\documentclass[11pt,letterpaper]{article}
\usepackage[T1]{fontenc}
\usepackage{lmodern,microtype}
\usepackage[margin=.75in,headheight=15pt]{geometry}
\usepackage{amsmath,amssymb,mathtools}
\usepackage{xcolor,enumitem,fancyhdr,hyperref}
\definecolor{ink}{HTML}{222834}
\definecolor{indigo}{HTML}{25355D}
\definecolor{muted}{HTML}{505A6B}
\definecolor{keyink}{HTML}{205D52}
\color{ink}
\hypersetup{colorlinks=true,linkcolor=indigo,urlcolor=indigo,
  pdftitle={Probability Midterm 1: Makeup},pdfauthor={Igor Rivin}}
\setlength{\parindent}{0pt}
\setlength{\parskip}{4pt}
\setlength{\abovedisplayskip}{6pt}
\setlength{\belowdisplayskip}{6pt}
\clubpenalty=10000
\widowpenalty=10000
\setlist[enumerate]{label=(\alph*),leftmargin=1.7em,
  itemsep=4pt,topsep=5pt,parsep=2pt}
\newcommand{\Prob}{\mathbb P}
\newcommand{\E}{\mathbb E}
\DeclareMathOperator{\Var}{Var}
\newcommand{\edition}{\ifanswers Instructor key\else Student paper\fi}
\newcommand{\partpts}[1]{\unskip\nobreak\hspace{.3em}%
  \mbox{\small\color{muted}(#1 points)}}
\newcommand{\credit}[1]{\par{\footnotesize\color{muted}\textit{Credit:} #1}}
\newcommand{\problem}[3]{\par\addvspace{9pt}%
  {\large\bfseries\color{indigo}#1.\ #2}\hfill
  {\small\color{muted}\textbf{#3 points}}\par\nobreak}
\NewDocumentEnvironment{solution}{O{1in}+b}{%
  \ifanswers
    \par\begingroup\small\color{keyink}\textbf{Solution.} #2\par\endgroup
  \else
    \par\nobreak\vspace*{#1}
  \fi
}{}
\pagestyle{fancy}
\fancyhf{}
\fancyhead[L]{\small\color{muted}Probability \textbar\ Midterm 1 Makeup}
\fancyhead[R]{\small\color{muted}Chance, Data, and Decisions}
\fancyfoot[L]{\footnotesize\color{muted}\edition}
\fancyfoot[R]{\footnotesize\color{muted}Page \thepage\ of \pageref*{exam-last-page}}
\renewcommand{\headrulewidth}{.3pt}
\begin{document}
{\LARGE\bfseries\color{indigo}Probability Midterm 1: Makeup}\par
{\small\color{muted}50 minutes \quad $\bullet$ \quad 100 points
  \quad $\bullet$ \quad Six problems}\par
\ifanswers
\textbf{Grading note.} Award method credit after an earlier error; do not
deduct twice for the same mistake. Suggested working times for Problems
1--6 are 5, 5, 6, 8, 10, and 9 minutes, leaving 7 minutes to check.
\else
\smallskip
\textbf{Name:}\ \rule{.49\linewidth}{.4pt}\hfill
\textbf{Date:}\ \rule{.20\linewidth}{.4pt}\par
\fi
\textbf{Instructions.} Answer all six problems and show your reasoning.
Exact answers are acceptable; binomial coefficients need not be evaluated.
No calculator is needed. Use the spaces provided; clearly label any work
continued on the back. In Problem 6, use the guided steps rather than
quoting a coupon-collector formula. Reserve time to check your work.

\problem{1}{Sets}{10}
Let $U=\{1,2,\ldots,12\}$, $A=\{1,3,5,7,9,11\}$, and
$B=\{1,2,3,4,5\}$. Take all complements relative to $U$.
\begin{enumerate}
\item Find $A\cap B$, $A\setminus B$, and $(A\cup B)^c$. \partpts{6}
\begin{solution}[.85in]
\[
A\cap B=\{1,3,5\},\quad A\setminus B=\{7,9,11\},\quad
(A\cup B)^c=\{6,8,10,12\}.
\]
\credit{2 per correct set; allow partial credit for a single entry error.}
\end{solution}
\item For these sets $A,B$, is the following statement true for every
$C\subseteq U$? Justify with a proof or a counterexample. \partpts{4}
\[
C\subseteq A\cup B\quad\Longrightarrow\quad
C\subseteq A\ \text{or}\ C\subseteq B.
\]
\begin{solution}[.85in]
False: take $C=\{2,7\}$. Both elements belong to $A\cup B$, but
$2\notin A$ and $7\notin B$. Hence neither containment on the right holds.
\credit{False: 1; valid example: 1; union containment: 1;
both failed containments: 1. Accept any valid counterexample.}
\end{solution}
\end{enumerate}

\problem{2}{A probability space}{10}
Let $\Omega=\{a,b,c,d\}$, with every subset an event. The probabilities
of $a,b,c,d$ are $2k,3k,4k,5k$, respectively.
\begin{enumerate}
\item Find $k$, explaining the probability axiom you use. \partpts{4}
\begin{solution}[.6in]
The disjoint singletons partition $\Omega$. Additivity and
$\Prob(\Omega)=1$ give $14k=1$, so $k=1/14$. The probabilities are nonnegative.
\credit{Normalization/additivity explanation: 2; equation: 1; value: 1.}
\end{solution}
\item Let $E=\{a,c\}$ and $F=\{b,c\}$. Find $\Prob(E)$,
$\Prob(E\cap F)$, and $\Prob(E^c)$. \partpts{6}
\begin{solution}[.85in]
Since $E\cap F=\{c\}$ and $E^c=\{b,d\}$,
\[
\Prob(E)=\frac6{14}=\frac37,\qquad
\Prob(E\cap F)=\frac4{14}=\frac27,\qquad
\Prob(E^c)=\frac8{14}=\frac47.
\]
\credit{2 per probability. Carry forward an earlier error in $k$:
accept the correctly obtained expressions $6k,4k,8k$.}
\end{solution}
\end{enumerate}

\newpage
\problem{3}{Inclusion--exclusion}{15}
In a group of 45 students, 20 play tennis, 17 play basketball, and 13 play
volleyball. Eight play both tennis and basketball, six both tennis and
volleyball, and five both basketball and volleyball. Three play all three.
\textbf{The pairwise counts include students who play all three sports.}
\begin{enumerate}
\item How many students play at least one of these sports? Use
inclusion--exclusion. \partpts{11}
\begin{solution}[.95in]
Writing $T,B,V$ for the three sets,
\begin{align*}
|T\cup B\cup V|
&=|T|+|B|+|V|-|T\cap B|-|T\cap V|-|B\cap V|
 +|T\cap B\cap V|\\
&=20+17+13-8-6-5+3=34.
\end{align*}
\credit{Individual counts: 3; subtracting pairwise intersections: 3;
restoring the triple intersection: 3; result: 2.}
\end{solution}
\item How many students play none of these sports? \partpts{4}
\begin{solution}[.65in]
There are $45-34=11$ such students.
\credit{Complement subtraction: 3; result: 1. Carry forward an earlier error.}
\end{solution}
\end{enumerate}

\problem{4}{Balls in an urn}{20}
An urn contains six red, three blue, and four green balls. Four balls are
selected uniformly without replacement: every four-ball subset of the
thirteen individual balls is equally likely.
\begin{enumerate}
\item Find the probability of selecting exactly two red, one blue,
and one green ball. \partpts{10}
\begin{solution}[1.25in]
Choose the required number of each color and divide by the total number
of four-ball subsets:
\[
\Prob(\text{two red, one blue, one green})
=\frac{\binom62\binom31\binom41}{\binom{13}4}
=\frac{180}{715}=\frac{36}{143}.
\]
\credit{Uniform denominator: 3; red choice: 2; blue and green choices: 2;
product and ratio: 3. No numerical simplification needed.}
\end{solution}
\item Find the probability of selecting at least one green ball. \partpts{10}
\begin{solution}[1.25in]
The nine non-green balls give the complementary count:
\[
\Prob(\text{at least one green})
=1-\frac{\binom94}{\binom{13}4}
=1-\frac{126}{715}=\frac{589}{715}.
\]
\credit{Complement event: 3; favorable complement count: 3;
denominator: 2; subtracting from 1: 2. Accept an equivalent direct count.}
\end{solution}
\end{enumerate}

\newpage
\problem{5}{A density, its mean, and its variance}{25}
A continuous random variable $X$ has probability density
\[
f(x)=\begin{cases}cx^2,&0\le x\le1,\\0,&\text{otherwise},\end{cases}
\]
where $c$ is a constant. Show the integrals you use.
\begin{enumerate}
\item Find $c$. \partpts{5}
\begin{solution}[1.4in]
Normalization requires
\[
1=c\int_0^1x^2\,dx=c\left[\frac{x^3}{3}\right]_0^1=\frac c3.
\]
Thus $c=3$, which also makes the density nonnegative.
\credit{Normalization integral: 2; integral value: 2; constant: 1.}
\end{solution}
\item Compute $\E[X]$. \partpts{8}
\begin{solution}[1.5in]
\[
\E[X]=\int_0^1x(3x^2)\,dx
=3\left[\frac{x^4}{4}\right]_0^1=\frac34.
\]
\credit{Moment integral including the factor $x$: 4;
antiderivative and limits: 3; result: 1. Carry forward an earlier value of $c$.}
\end{solution}
\item Compute $\E[X^2]$, and hence $\Var(X)$. \partpts{12}
\begin{solution}[2in]
\[
\E[X^2]=\int_0^1x^2(3x^2)\,dx
=3\left[\frac{x^5}{5}\right]_0^1=\frac35.
\]
Therefore
\[
\Var(X)=\E[X^2]-(\E[X])^2
=\frac35-\left(\frac34\right)^2
=\frac{48-45}{80}=\frac3{80}.
\]
\credit{Second-moment integral: 4; evaluation: 3; variance identity: 3;
substitution and arithmetic: 2. Accept a correct direct variance integral.}
\end{solution}
\end{enumerate}

\newpage
\problem{6}{Collecting five sticker types}{20}
Each packet contains one sticker. There are five sticker types, each equally
likely, and different packets have independent contents. Duplicates are
possible. Starting with no stickers, let $T$ be the number of packets opened
up to and including the first packet that completes the collection of all
five types.
\begin{enumerate}
\item After exactly $k$ distinct types have been collected, let $p_k$ be
the probability that the next packet contains a new type. Find
$p_0,p_1,p_2,p_3,p_4$, and explain your reasoning. \partpts{6}
\begin{solution}[1.3in]
There are $5-k$ missing types out of five equally likely types, so
\[
p_k=\frac{5-k}{5},\qquad
(p_0,p_1,p_2,p_3,p_4)=\left(1,\frac45,\frac35,\frac25,\frac15\right).
\]
\credit{Missing-type counting argument: 1; five correct probabilities: 1 each.}
\end{solution}
\item Let $W_k$ be the number of additional packets needed to go from
exactly $k$ collected types to $k+1$ collected types, including the packet
containing the new type. Find $\E[W_k]$ for $k=0,1,2,3,4$. \partpts{8}

\textit{You may use:} A geometric waiting time that counts trials through
the first success, with success probability $p$, has mean $1/p$.
\begin{solution}[1.6in]
Until a new type appears, the missing set stays fixed. Independent packets
therefore give a geometric wait with success probability $p_k$ (the first
wait is deterministically one). Thus
\[
\bigl(\E[W_0],\E[W_1],\E[W_2],\E[W_3],\E[W_4]\bigr)
=\left(1,\frac54,\frac53,\frac52,5\right).
\]
\credit{Applying the geometric reciprocal rule: 3; five means: 1 each.
Carry forward earlier probabilities when reciprocals are taken consistently.}
\end{solution}
\item Write $T$ as a sum of the waiting times and find $\E[T]$.
Name the expectation rule used. \partpts{6}
\begin{solution}[1.8in]
The stages partition all packets opened, so $T=\sum_{k=0}^4W_k$.
Linearity of expectation gives
\[
\E[T]=1+\frac54+\frac53+\frac52+5=\frac{137}{12}.
\]
Independence of the waiting times is not required for linearity.
\credit{Stage decomposition: 2; linearity and sum: 2; evaluation: 2.
Accept the exact unevaluated sum.}
\end{solution}
\end{enumerate}
\vfill
\begin{center}\small\color{muted}End of paper. Check that you have answered all six problems.\end{center}
\label{exam-last-page}
\end{document}
